% FPL 2026 Short Paper — 4 pages + 1 references
% IEEE Conference format
\documentclass[conference]{IEEEtran}

\usepackage{cite}
\usepackage{amsmath,amssymb}
\usepackage{graphicx}
\usepackage{textcomp}
\usepackage{booktabs}
\usepackage{xcolor}
\usepackage{url}
\usepackage{multirow}

\newcommand{\tmark}[1]{\textsuperscript{#1}}

\begin{document}

\title{Zero-DSP Ternary Transformer Inference\\on a \$30 FPGA with Open-Source Toolchain}

\author{
\IEEEauthorblockN{Dmitrii Vasilev}
\IEEEauthorblockA{Trinity Project\\
\texttt{https://github.com/gHashTag/trinity}}
}

\maketitle

% ====================================================================
\begin{abstract}
We present Trinity, a ternary transformer inference engine on a \$30 Xilinx
Artix-7 FPGA (XC7A100T) using zero DSP48 blocks and a fully open-source
synthesis toolchain (Yosys, nextpnr-xilinx, Project X-Ray).  Weights are
constrained to $\{-1, 0, +1\}$, reducing multiply-accumulate to pure
add/subtract logic.

The core contribution is the Ternary MatMul Unit (TMU): a $K\!=\!32$ parallel
dot-product engine that reads 32 ternary weights per cycle from banked BRAMs
and reduces them through a 5-stage combinational adder tree, achieving
$26\times$ speedup over sequential accumulation.

A 3-block autoregressive pipeline with full TMU coverage (including LM~head)
generates tokens at a simulation-verified rate of ${\sim}1{,}879$~tok/s at
81.25~MHz.  The design consumes ${\sim}20$K~LUTs (32\%), ${\sim}117.5$
BRAM36-equivalent tiles (87\%), and \textbf{0~DSP48E1} blocks, storing
1,125,090 ternary weights on-chip.

Compared to TerEffic (U280, \$17K, 2,733 DSPs), TeLLMe~v2 (KV260, \$249,
1,200 DSPs), and FlightLLM (U280, 2,733 DSPs), Trinity achieves the lowest
cost per tok/s at \$0.016.  While our model (1.95M parameters) is much smaller
than prior work (125M--7B), the result demonstrates that ternary transformers
can fully saturate low-cost FPGAs without DSP blocks, opening a new cost regime
for edge inference.  To our knowledge, this is the first zero-DSP ternary
transformer on a sub-\$50 FPGA with open-source synthesis.
\end{abstract}

\begin{IEEEkeywords}
Ternary neural network, FPGA, zero-DSP, open-source EDA, edge inference
\end{IEEEkeywords}

% ====================================================================
\section{Introduction}

Deploying neural networks on edge devices is constrained by power, cost, and
silicon area.  GPU and high-end FPGA solutions achieve high throughput at
costs of \$5,000--\$17,000 and power budgets of 27--40\,W, precluding
ultra-low-cost applications.

Ternary quantization---constraining weights to $\{-1, 0, +1\}$---eliminates
multiplications: each weight--input product reduces to addition, subtraction,
or no-op.  This makes ternary networks uniquely suited to low-cost FPGAs
where DSP resources are scarce.

Recent work targets high-end devices.  TerEffic~\cite{tereffic} achieves
16,300~tok/s on Alveo~U280 (\$17K) with 2,733~DSPs.  TeLLMe~v2~\cite{tellme}
targets KV260 (\$249) with 1,200~DSPs.  FlightLLM~\cite{flightllm} uses U280
for 153~tok/s.  LUT-LLM~\cite{lutllm} targets V80 (\$6K) with 2,880~DSPs.
All require Vivado.

Trinity differs in three dimensions:
\begin{enumerate}
  \item \textbf{Zero DSP.}  The entire datapath---matrix-vector multiplication,
    ReLU, residual, RMSNorm---uses purely LUT-based logic.
  \item \textbf{Open-source toolchain.}  Yosys~0.63~\cite{yosys} (synthesis),
    nextpnr-xilinx~\cite{nextpnr} (P\&R), Project~X-Ray~\cite{prjxray}
    (bitstream) require no licenses.
  \item \textbf{Sub-\$50 target.}  The XC7A100T on a QMTECH board costs
    ${\sim}\$30$, two orders of magnitude below datacenter FPGAs.
\end{enumerate}

The key architectural contribution is the \textit{Ternary MatMul Unit}~(TMU),
a $K\!=\!32$ parallel dot-product engine using banked BRAMs and a 5-stage
combinational adder tree.

% ====================================================================
\section{Architecture}

\subsection{System Overview}

The inference engine implements autoregressive token generation:
\[
\text{token\_id} \to \text{Emb} \to B_1 \to B_2 \to B_3
\to \text{LM\_Head} \to \text{Argmax} \to \text{next\_token}
\]
Three TrinityBlocks process a $3^5\!=\!243$-dimensional vector through
up-projection to $3^6\!=\!729$ dimensions, ReLU activation, down-projection,
residual connection, and shift-based RMSNorm.  Blocks execute sequentially;
inter-block data passes through distributed RAM buffers.

\subsection{Ternary MatMul Unit (TMU)}

The TMU replaces sequential weight-by-weight accumulation with $K\!=\!32$
parallel processing.

\textbf{Weight banking.}  Each matrix $\mathbf{W} \in \{-1,0,+1\}^{N_\text{in}
\times N_\text{out}}$ is distributed across 32 BRAM banks.  Bank~$b$ stores
weights where column index $i \bmod 32 = b$, at address $j \cdot
\lceil N_\text{in}/32 \rceil + \lfloor i/32 \rfloor$.  All banks share a
single read address, delivering 32 weight codes per cycle.

\textbf{Ternary mux.}  Each 2-bit code selects the operation:
\texttt{01}$\to$~$+x_b$, \texttt{10}$\to$~$-x_b$, \texttt{00}$\to$~$0$.
No multiplier is used.

\textbf{Adder tree.}  A 5-stage combinational tree ($32\to 16\to 8\to 4\to
2\to 1$) sums the 32 partial products.  The output accumulates across
$\lceil N_\text{in}/32 \rceil$ steps per output element.

\textbf{Cycle count.}  For $243\!\to\!729$ (up): $\lceil 243/32\rceil = 8$
steps/output $\times$ 729~outputs $\approx$ 6.8K~cycles.  For
$729\!\to\!243$ (down): 23~steps $\times$ 243~outputs $\approx$ 6.1K~cycles.
Per block: ${\sim}13.9$K~cycles vs ${\sim}360$K for $K\!=\!1$
($26\times$ speedup).

\subsection{TrinityBlock}

Each block computes:
\begin{align}
  \mathbf{h} &= \text{TMU}_\text{up}(\mathbf{x}) & 243 &\to 729 \\
  \mathbf{a} &= \text{ReLU}(\mathbf{h}) \\
  \mathbf{d} &= \text{TMU}_\text{down}(\mathbf{a}) & 729 &\to 243 \\
  \mathbf{r} &= \mathbf{d} + \mathbf{x} & \multicolumn{2}{l}{\text{(residual)}} \\
  \mathbf{y} &= \text{ShiftRMSNorm}(\mathbf{r})
\end{align}

\textbf{Shift-based RMSNorm} replaces division with priority-encoder MSB
detection and barrel-shift normalization, using only LUT logic.

\textbf{BRAM budget.}  Each TMU uses 32~BRAM18 banks (2-bit $\times$ depth),
yielding 16~BRAM36-equivalent per TMU.  Per block: $2 \times 16 = 32$~BRAM36-eq.
Three blocks: 96~BRAM36-eq.  The LM~head TMU ($243\!\to\!128$) adds 16~BRAM36-eq.
With embedding: ${\sim}117.5$~BRAM36-eq total (87\% of device capacity).

\subsection{Self-Test}

An autonomous power-on test feeds seed token~42 through the full pipeline,
generating 16~tokens autoregressively.  Block~3 output is validated
(243~elements, $\ge 1$ non-zero).  LED indicates PASS/FAIL; the token sequence
is reported via UART.

% ====================================================================
\section{Results}

\subsection{Throughput Evolution}

Table~\ref{tab:evolution} traces the four optimization steps from baseline
to the final design, each independently verified in simulation.

\begin{table}[t]
\centering
\caption{Optimization history: from 469 to 1,879 tok/s in four steps}
\label{tab:evolution}
\begin{tabular}{@{}lccrc@{}}
\toprule
\textbf{Step} & \textbf{$K$} & \textbf{Freq} & \textbf{Cyc/tok} & \textbf{tok/s} \\
\midrule
Baseline ($K\!=\!1$ LM head) & 16/1 & 50 & ${\sim}113$K & 469 \\
LM head $\to$ TMU $K\!=\!16$ & 16 & 50 & ${\sim}81.6$K & 613 \\
MMCM 50$\to$65~MHz            & 16 & 65 & 78,528 & 827 \\
TMU $K\!=\!16 \to 32$         & 32 & 65 & 43,241 & 1,503 \\
\textbf{MMCM 65$\to$81.25~MHz} & \textbf{32} & \textbf{81.25} & \textbf{43,241} & \textbf{1,879} \\
\bottomrule
\end{tabular}
\end{table}

The dominant gain comes from doubling TMU parallelism ($K\!=\!16 \to 32$):
block cycles drop from 25,286 to 13,865 ($-45\%$), LM~head from 2,422
to 1,398 ($-42\%$).  Raising the MMCM clock from 65 to 81.25~MHz contributes
an additional 25\% at zero cycle cost.

\textbf{Cycle breakdown} (simulation-verified, per token):
embedding 248, block$_1$ 13,865, block$_2$ 13,865, block$_3$ 13,865,
LM~head 1,398, argmax ${\sim}1$.  Total: 43,241~cycles.

\subsection{Resource Utilization}

Table~\ref{tab:resources} shows projected utilization for the $K\!=\!32$
design (Yosys synthesis pending; $K\!=\!16$ synthesis confirmed these
estimates track within 10\%).

\begin{table}[t]
\centering
\caption{Resource utilization: HSLM 3-block TMU $K\!=\!32$ pipeline}
\label{tab:resources}
\begin{tabular}{@{}lrrr@{}}
\toprule
\textbf{Resource} & \textbf{Est.\ used} & \textbf{Available} & \textbf{\%} \\
\midrule
LUT (INV+LUT2--6) & ${\sim}20{,}300$ & 63,400 & 32.1 \\
FF (FDRE+FDSE)     & ${\sim}2{,}500$  & 126,800 & 2.0 \\
BRAM36-eq          & ${\sim}117.5$ & 135     & 87.0 \\
DSP48E1            & \textbf{0} & 240 & \textbf{0} \\
\bottomrule
\end{tabular}
\end{table}

The design remains BRAM-limited at 87\% utilization.  LUT usage grows
modestly ($+12\%$) from the wider adder tree (16 additional 20-bit adders
at stage~1).

\subsection{Comparison with Prior Work}

Table~\ref{tab:comparison} compares Trinity with recent FPGA ternary
accelerators.

\begin{table}[t]
\centering
\caption{Comparison with FPGA ternary inference engines}
\label{tab:comparison}
\setlength{\tabcolsep}{3pt}
\begin{tabular}{@{}lcrrrl@{}}
\toprule
\textbf{System} & \textbf{FPGA} & \textbf{Cost} & \textbf{DSP} & \textbf{tok/s} & \textbf{Tools} \\
\midrule
\textbf{Trinity} & A7-100T & \textbf{\$30} & \textbf{0} & \textbf{1,879\tmark{*}} & \textbf{Open} \\
TeLLMe~v2 & KV260 & \$249 & 1,200 & 25 & Vivado \\
FlightLLM & U280 & \$17K & 2,733 & 153 & Vivado \\
LUT-LLM & V80 & \$6K & 2,880 & 175 & Vivado \\
TerEffic & U280 & \$17K & 2,733 & 16,300 & Vivado \\
\bottomrule
\multicolumn{6}{l}{\footnotesize \tmark{*}Simulation-verified at 81.25~MHz; hardware timing pending.}
\end{tabular}
\end{table}

\textbf{Cost efficiency.}  \$30\,/\,1{,}879~tok/s $= \$0.016$/tok/s, the
lowest in Table~\ref{tab:comparison}.  TeLLMe: \$9.96; FlightLLM: \$113;
TerEffic: \$1.07.

\textbf{Timing note.}  We measure 43,241 cycles/token in simulation, yielding
1,879~tok/s at 81.25~MHz.  The 5-stage adder tree critical path
(${\sim}10$~ns) fits within the 12.3~ns clock period.  Timing closure at this
frequency awaits static timing analysis with place-and-route.

\subsection{Limitations}

We state the following honestly:
\begin{enumerate}
  \item Throughput (${\sim}1{,}879$~tok/s) is derived from cycle counts at
    81.25~MHz.  Hardware UART timing measurement is pending.
  \item Trained weights are \emph{not} loaded on FPGA.  The bitstream uses
    random ternary weights.  Self-test validates the datapath, not model
    quality.
  \item Power is not measured.  We make no tok/s/W claims.  Power and energy
    efficiency measurements are left as future work, pending hardware kit
    arrival in late March~2026.
  \item 1.95M ternary parameters is proof-of-concept scale.
    This is not competitive with BitNet~2.4B on text quality.
  \item Attention is not implemented on FPGA (BRAM-limited on A7-100T).
  \item Resource estimates for $K\!=\!32$ await Yosys synthesis confirmation
    ($K\!=\!16$ synthesis is verified).
\end{enumerate}

% ====================================================================
\section{Conclusion}

We presented Trinity, a 3-block ternary transformer on a \$30 FPGA using zero
DSP blocks and open-source tools.  Through four incremental optimizations---TMU
parallelism ($K\!=\!1 \to 16 \to 32$) and clock frequency (50$\to$81.25~MHz)---
we improve throughput from 469 to 1,879~tok/s, a $4.0\times$ gain at constant
silicon cost.

The final design processes 43,241 cycles/token using ${\sim}20$K~LUTs,
${\sim}117.5$~BRAM36-eq, and 0~DSP48, achieving \$0.016/tok/s---the lowest
reported cost among FPGA ternary inference engines.  The zero-DSP TMU
architecture scales independently of model quality: larger models on bigger
FPGAs would inherit the same cost advantage.

Scaling paths include XC7A200T (365~BRAM36, ${\sim}9$ blocks), Kintex-7
(445~BRAM36, 200+~MHz), and eFPGA integration.

The design is fully open-source:
\url{https://github.com/gHashTag/trinity}

% ====================================================================
\bibliographystyle{IEEEtran}
\begin{thebibliography}{8}

\bibitem{tereffic}
C.~Yin \textit{et al.}, ``TerEffic: Highly Efficient Ternary LLM Inference
on FPGA,'' \textit{arXiv:2502.16473v2}, May 2025.

\bibitem{tellme}
H.~Li \textit{et al.}, ``TeLLMe: An End-to-End Framework for LLM Deployment
on Low-Power FPGAs,'' 2025.

\bibitem{flightllm}
S.~Lu \textit{et al.}, ``FlightLLM: Efficient Large Language Model Inference
with a Complete Mapping Flow on FPGAs,'' \textit{FPGA}, 2024.

\bibitem{lutllm}
W.~Zhang \textit{et al.}, ``LUT-LLM: LUT-Based Efficient LLM Inference on
FPGA,'' 2025.

\bibitem{twn}
M.~Courbariaux \textit{et al.}, ``Ternary Weight Networks,''
\textit{arXiv:1605.04711}, 2016.

\bibitem{yosys}
C.~Wolf, \textit{Yosys Open SYnthesis Suite},
\url{https://github.com/YosysHQ/yosys}.

\bibitem{nextpnr}
\textit{nextpnr-xilinx},
\url{https://github.com/openXC7/nextpnr-xilinx}.

\bibitem{prjxray}
\textit{Project X-Ray},
\url{https://github.com/f4pga/prjxray}.

\end{thebibliography}

\end{document}
